\subsection{Propiedades}

\subsubsection{Regla: PascalCase y acceso mínimo necesario}

\textbf{Regla:} Las propiedades se nombran en \texttt{PascalCase} con un sustantivo, frase nominal o adjetivo. Si su valor no cambia después de construir el objeto, se expone solo \texttt{get} o se utiliza \texttt{init}; no se agrega un \texttt{set} público innecesario (Microsoft, 2025f).

\textbf{Código correcto:}
\begin{lstlisting}[style=csharpstyle]
public sealed class PlayerProfile
{
    public PlayerProfile(string displayName)
    {
        DisplayName = displayName;
    }

    public string DisplayName { get; }
}
\end{lstlisting}

\textbf{Código incorrecto:}
\begin{lstlisting}[style=csharpstyle]
public sealed class PlayerProfile
{
    public string displayName { get; set; }
}
\end{lstlisting}

\subsubsection{Regla: modificación controlada con \texttt{private set}}

\textbf{Regla:} Cuando una propiedad cambia durante la partida, se mantiene el \texttt{get} público y se restringe el \texttt{set} al menor nivel necesario. La modificación del estado de dominio se realiza mediante un método que expresa la operación (Microsoft, 2025f).

\textbf{Código correcto:}
\begin{lstlisting}[style=csharpstyle]
public sealed class Player
{
    public int Score { get; private set; }

    public void AddScore(int points)
    {
        Score += points;
    }
}
\end{lstlisting}

\textbf{Código incorrecto:}
\begin{lstlisting}[style=csharpstyle]
public sealed class Player
{
    public int Score { get; set; }
}
\end{lstlisting}
